\documentclass[UTF8,zihao=-4]{ctexart}

\usepackage[a4paper,margin=2.4cm]{geometry}
\usepackage{amsmath,amssymb,bm}
\usepackage{booktabs}
\usepackage{array}
\usepackage{graphicx}
\usepackage{xcolor}

\title{\textbf{面向SAR舰船检测的去斑点选择性状态空间建模方法框架}}
\author{Method Draft for DS-SSM / IG-DS-SSM / DG-DS-SSM}
\date{\today}

\newcommand{\x}{\mathbf{x}}
\newcommand{\xtilde}{\widetilde{\mathbf{x}}}
\newcommand{\h}{\mathbf{h}}
\newcommand{\I}{\mathbf{I}}
\newcommand{\R}{\mathbf{R}}

\begin{document}
\maketitle

\begin{abstract}
SAR舰船检测面临显著的斑点噪声、海杂波干扰以及小目标弱纹理表达问题。直接将通用视觉骨干或Mamba类状态空间模块用于SAR图像时，模型可能将噪声扰动与真实舰船结构一并建模，导致状态传播过程中的目标响应被背景噪声污染。为此，本文方法部分建议围绕De-Speckle Selective State Space Model（DS-SSM）展开：首先通过局部背景估计构建去斑点输入表征，再利用能量感知抑制机制约束状态传播；随后引入Perceptual Feature Importance Module（PFIM）估计目标重要性，并将其注入记忆保留、信息写入和动态步长调制过程。作为主线方法的延伸，进一步探索Decoupled-Gated DS-SSM（DG-DS-SSM），将状态更新显式拆分为擦除、写入和步长调制三个分支，以获得更细粒度的SAR目标与噪声状态控制。
\end{abstract}

\tableofcontents
\newpage

\section{方法部分整体组织建议}

按照顶会论文的写法，方法部分不宜只罗列模块，而应围绕一个清晰问题展开：\textbf{如何在状态空间建模中抑制SAR斑点噪声，同时保留舰船目标的有效结构响应}。建议方法章节组织如下：

\begin{enumerate}
  \item \textbf{Overall Architecture}：说明本文方法嵌入Mamba-YOLO检测框架，用DS-SSM替换或增强原有Mamba视觉块。
  \item \textbf{De-Speckle Selective State Space Model}：介绍DS-SSM核心，包括局部背景估计、输入去偏置、能量感知状态抑制。
  \item \textbf{Perceptual Feature Importance Guided Modulation}：介绍PFIM如何估计重要性图，并调制状态保留、写入和步长。
  \item \textbf{Decoupled-Gated Extension}：作为延长线介绍DG-DS-SSM，即erase/write/delta解耦门控。
  \item \textbf{Training Strategy and Inference}：介绍20 epoch warmup、单卡/分布式训练一致性、SSDD上指标选择。
\end{enumerate}

\section{Overall Architecture}

给定输入SAR图像$\mathbf{X}\in\mathbb{R}^{H\times W\times C}$，检测骨干首先提取多尺度特征。本文方法将常规视觉状态空间块替换为SAR感知的DS-SSM块。对于第$l$层特征$\mathbf{F}^{l}$，首先将二维特征展平成token序列：
\begin{equation}
  \mathbf{x} = \mathrm{Flatten}(\mathbf{F}^{l}) \in \mathbb{R}^{N\times C}, \quad N=H_l W_l .
\end{equation}
随后，DS-SSM在token序列上执行去斑点选择性状态建模：
\begin{equation}
  \mathbf{y} = \mathrm{DS\mbox{-}SSM}(\mathbf{x}),
\end{equation}
并通过残差连接和前馈网络得到输出：
\begin{equation}
  \mathbf{z} = \mathbf{x} + \gamma_1 \mathrm{DS\mbox{-}SSM}(\mathrm{LN}(\mathbf{x})), \quad
  \mathbf{o} = \mathbf{z} + \gamma_2 \mathrm{FFN}(\mathrm{LN}(\mathbf{z})).
\end{equation}

该设计保持Mamba-YOLO的高效序列建模能力，同时将SAR图像中的斑点噪声抑制机制显式注入状态空间更新过程。

\section{De-Speckle Selective State Space Model}

\subsection{Motivation}

SAR图像中的斑点噪声通常表现为局部高频扰动。若直接将输入特征送入状态空间递推，噪声token会沿序列传播并影响后续状态。DS-SSM的核心思想是：\textbf{在状态更新之前先构建去斑点输入，并在状态传播中根据局部能量抑制不可靠响应}。

\subsection{Local Background Estimation and Input De-biasing}

对输入token序列$\mathbf{x}$，首先通过深度可分离局部平滑估计局部背景：
\begin{equation}
  \bar{\mathbf{x}} = \mathrm{DWConv}(\mathbf{x}).
\end{equation}
局部残差为：
\begin{equation}
  \mathbf{r} = \mathbf{x} - \bar{\mathbf{x}} .
\end{equation}
DS-SSM通过可学习系数$\lambda_s$抑制局部斑点残差，得到去斑点输入：
\begin{equation}
  \widetilde{\mathbf{x}} = \mathbf{x} - \lambda_s \mathbf{r}.
\end{equation}
其中$\lambda_s$控制去斑点强度。较大的$\lambda_s$可以更强地抑制局部扰动，但也可能损失细粒度目标结构，因此实际实现中应对其范围进行约束。

\subsection{Energy-aware State Damping}

SAR背景海杂波区域通常具有不稳定局部能量。本文利用局部背景能量估计token可靠性：
\begin{equation}
  e_t = \frac{\|\bar{\mathbf{x}}_t\|_2}{\sqrt{C}} .
\end{equation}
随后构建状态保留门：
\begin{equation}
  g_t = g_{\min} + (1-g_{\min})\sigma(-\alpha \sigma_n e_t),
\end{equation}
其中$g_{\min}$为最小保留率，$\sigma_n$表示噪声强度先验，$\alpha$为可学习调制系数。该门控使高能量不稳定区域的状态传播受到抑制，从而降低斑点噪声沿状态空间递推扩散的风险。

\subsection{Selective State Update}

在去斑点输入$\widetilde{\mathbf{x}}$上生成SSM输入、门控和状态参数：
\begin{equation}
  \mathbf{u},\mathbf{z}=\mathrm{Linear}(\widetilde{\mathbf{x}}), \quad
  \Delta_t=\mathrm{softplus}(\mathrm{Linear}(\mathbf{u}_t)).
\end{equation}
状态递推可写为：
\begin{equation}
  \mathbf{h}_t = g_t \odot \bar{\mathbf{A}}_t \mathbf{h}_{t-1}
  + \bar{\mathbf{B}}_t \mathbf{u}_t,
\end{equation}
\begin{equation}
  \mathbf{y}_t = \mathbf{C}_t \mathbf{h}_t + \mathbf{D}\widetilde{\mathbf{x}}_t .
\end{equation}
其中$g_t$显式抑制噪声状态传播，$\mathbf{D}\widetilde{\mathbf{x}}_t$保留当前去斑点输入的残差信息。

\section{Perceptual Feature Importance Guided Modulation}

\subsection{PFIM Motivation}

仅依靠局部能量抑制可能误伤舰船目标，因为真实舰船区域同样可能具有高响应。为解决这一问题，本文引入PFIM估计token级重要性，使模型能够区分“目标高响应”和“噪声高响应”。

\subsection{Importance Estimation}

PFIM以去斑点输入$\widetilde{\mathbf{x}}$、局部背景$\bar{\mathbf{x}}$以及局部统计信息为输入，生成重要性图：
\begin{equation}
  \mathbf{I} = \mathrm{PFIM}(\widetilde{\mathbf{x}}, \bar{\mathbf{x}}, \mathrm{Std}_{local}(\mathbf{x})),
\end{equation}
其中$\mathbf{I}\in[0,1]^{N\times 1}$表示每个token的重要性。局部标准差为PFIM提供纹理和斑点强度线索，有助于识别弱目标和复杂海杂波。

\subsection{Importance-guided Retention, Writing and Transition}

PFIM的重要性图被注入三个状态更新环节。

\paragraph{Memory retention.}
重要性调制状态保留门：
\begin{equation}
  g_t = g_{\min} + (1-g_{\min})
  \sigma(-\alpha\sigma_n e_t + \beta(I_t-0.5)).
\end{equation}
重要目标区域获得更高保留率，噪声区域被更强抑制。

\paragraph{State writing.}
写入门控制当前输入写入状态的强度：
\begin{equation}
  w_t = 1 + \rho \tanh(I_t-0.5),
\end{equation}
\begin{equation}
  \mathbf{u}^{write}_t = w_t \odot \mathbf{u}_t.
\end{equation}

\paragraph{Dynamic transition.}
重要性进一步调制状态转移步长：
\begin{equation}
  \Delta_t = \mathrm{softplus}(\mathrm{Linear}(\mathbf{u}_t) + \eta I_t).
\end{equation}
该设计允许模型在目标区域采用更灵活的动态响应，在噪声区域保持更保守的状态更新。

\section{DG-DS-SSM: Decoupled-Gated Extension}

\subsection{Motivation}

DS-SSM和PFIM已经证明重要性感知状态更新对SSDD有效。为了进一步提升状态控制能力，本文将其作为延长线提出DG-DS-SSM。其核心思想是将原先耦合的状态更新拆分为三个更可解释的分支：erase gate、write gate和delta gate。

\subsection{Decoupled Erase Gate}

Erase gate负责控制旧状态的保留或擦除：
\begin{equation}
  e^{erase}_t = -\alpha\sigma_n e_t + \beta_e(I_t-0.5),
\end{equation}
\begin{equation}
  r_t = r_{\min} + (1-r_{\min})\sigma(e^{erase}_t).
\end{equation}
该分支用于抑制斑点噪声和海杂波在历史状态中的持续传播。

\subsection{Decoupled Write Gate}

Write gate负责控制当前输入写入状态的强度：
\begin{equation}
  q_t = 1 + s_w \tanh(\rho_w(I_t-0.5)),
\end{equation}
\begin{equation}
  \mathbf{u}^{write}_t = q_t \odot \mathbf{u}_t.
\end{equation}
当PFIM判断当前token更可能属于舰船目标时，模型允许更多有效结构写入状态。

\subsection{Delta Gate}

Delta gate控制状态转移步长：
\begin{equation}
  \Delta_t = \mathrm{clip}\left(
  \mathrm{softplus}(\mathrm{Linear}(\mathbf{u}_t)+\eta_d I_t)
  \right).
\end{equation}
该分支影响SSM的动态响应速度，使目标区域和背景区域具有不同的状态演化强度。

\subsection{Unified Update}

DG-DS-SSM的递推过程可写为：
\begin{equation}
  \mathbf{h}_t = r_t \odot \bar{\mathbf{A}}_t \mathbf{h}_{t-1}
  + \bar{\mathbf{B}}_t \mathbf{u}^{write}_t,
\end{equation}
\begin{equation}
  \mathbf{y}_t = \mathbf{C}_t \mathbf{h}_t + \mathbf{D}\widetilde{\mathbf{x}}_t.
\end{equation}
与DS-SSM相比，DG-DS-SSM将“保留旧状态”“写入新信息”“调整状态步长”显式解耦，因此具有更强的可解释性。

\paragraph{Forward procedure.}
The DS-SSM / DG-DS-SSM forward pass can be summarized as follows:
\begin{enumerate}
  \item Estimate local background $\bar{\mathbf{x}}=\mathrm{DWConv}(\mathbf{x})$.
  \item Compute residual $\mathbf{r}=\mathbf{x}-\bar{\mathbf{x}}$.
  \item Construct de-speckled input $\widetilde{\mathbf{x}}=\mathbf{x}-\lambda_s\mathbf{r}$.
  \item Estimate importance $\mathbf{I}=\mathrm{PFIM}(\widetilde{\mathbf{x}},\bar{\mathbf{x}},H,W)$.
  \item Compute erase/retain gate $r_t$, write gate $q_t$, and dynamic step $\Delta_t$.
  \item Run selective state-space scan with gated input and retained state.
  \item Return output tokens with residual projection.
\end{enumerate}

\section{Training Strategy}

\subsection{Warmup for Stable Optimization}

由于DS-SSM包含去斑点系数、状态保留门和重要性调制分支，训练早期若所有参数同时更新，容易导致状态递推不稳定。当前工程实现采用20 epoch warmup策略：前20个epoch优先稳定新增模块或冻结部分主干参数，随后解冻全模型进行联合优化。该策略更适合作为训练稳定性设计，而不建议作为论文的核心理论创新。

\subsection{Single-GPU and DDP Consistency}

为了保证实验可复现，normal和DDP仓库应共享相同模型结构，仅训练启动方式不同。单卡实验建议使用：
\begin{equation}
  \texttt{CUDA\_VISIBLE\_DEVICES=k, --device 0}.
\end{equation}
分布式训练仅作为加速工具，不应改变模型模块定义。

\section{Experimental Positioning on SSDD}

SSDD是SAR舰船检测数据集，目标小且边界受SAR成像机理影响明显。因此，实验报告中建议以Precision、Recall和mAP50作为主要指标，mAP50-95作为辅助指标。当前关键结果可概括如下：

\begin{table}[h]
\centering
\caption{Current representative SSDD results.}
\begin{tabular}{lcccc}
\toprule
Method & Precision & Recall & mAP50 & mAP50-95 \\
\midrule
PFIM ex27 mainline & 0.9738 & 0.9522 & 0.9870 & 0.7402 \\
DG-DS-SSM ex34 standard & 0.9700 & 0.9560 & 0.9850 & 0.7350 \\
DG-DS-SSM ex34 TTA IoU0.55 & 0.9630 & 0.9620 & 0.9870 & 0.7240 \\
DG-DS-SSM ex32 write & 0.9540 & 0.9580 & 0.9860 & 0.7270 \\
\bottomrule
\end{tabular}
\end{table}

这些结果表明，DS-SSM/PFIM仍是当前最稳健的论文主线；DG-DS-SSM作为延长线在Recall和mAP50上具有竞争力，说明解耦门控状态更新能够增强SAR舰船目标检出能力。

\section{Recommended Writing Position}

论文中建议采用如下定位：

\begin{quote}
本文的核心贡献是提出面向SAR舰船检测的DS-SSM，通过去斑点输入修正、能量感知状态抑制和PFIM重要性引导，将SAR噪声先验注入选择性状态空间建模过程。进一步地，本文探索DG-DS-SSM作为DS-SSM的延伸，通过解耦erase、write和delta门控，使状态更新过程具备更细粒度的可解释调控能力。
\end{quote}

该写法能够突出DS-SSM作为核心创新，同时合理容纳今天新增的DG-DS-SSM实验，不会因为个别指标未全面超过ex27而削弱论文主线。

\section{Possible Contribution Statements}

\begin{enumerate}
  \item We propose DS-SSM, a SAR-oriented selective state-space module that suppresses speckle-induced state propagation through de-speckled input modeling and energy-aware damping.
  \item We introduce PFIM to estimate token-level target importance and inject it into memory retention, state writing and transition-step modulation.
  \item We further explore DG-DS-SSM, a decoupled-gated extension that explicitly separates erase, write and delta control for more interpretable SAR state updating.
  \item Extensive SSDD experiments show that the proposed design improves target detection reliability, especially in Recall and mAP50, while maintaining competitive localization performance.
\end{enumerate}

\end{document}
